\section*{Übersicht verwendeter Hilfsmittel}
\addcontentsline{toc}{section}{Übersicht verwendeter Hilfsmittel}
\label{app:hilfsmittel}

In dieser Arbeit wurden generative KI-Systeme als Hilfsmittel eingesetzt.

\begin{tabular}{@{}p{3.2cm}p{2.4cm}p{8cm}@{}}
\toprule
\textbf{System / Version} & \textbf{Zeitraum} & \textbf{Verwendung} \\
\midrule
Tabellen 21.8.2026 & 21.08.2026 & Erzeugung der Caption- und Notes-Texte aller
  37 Tabellen und der Abbildung der Arbeit (Prompt siehe unten) \\
PRODUKT X.Y & MM/JJJJ--MM/JJJJ & Entwurfsfassungen von Fließtext in
  \cref{sec:model-setup} und \cref{sec:conformal}, von mir überarbeitet \\
PRODUKT X.Y & MM/JJJJ--MM/JJJJ & Implementierung der Auswertungs-Pipeline
  (Python), \LaTeX-Satz \\
PRODUKT X.Y & MM/JJJJ--MM/JJJJ & Rechtschreib- und Grammatikkorrektur \\
\bottomrule
\end{tabular}

\bigskip
\noindent
Konzeption, Modellauswahl, Implementierung, Durchführung der Experimente und
Interpretation der Ergebnisse liegen ausschließlich bei mir.

\subsection*{Prompt für die Tabellenbeschreibungen}

Pro Tabelle einmal ausgeführt, mit dem Fließtext des Abschnitts, dem
Tabellencode und der Ergebnisdatei als Kontext. Der Prompttext war in allen
Durchläufen identisch. Der folgende Block gibt ihn in ASCII-Umschrift wieder;
das Zeichen $\dagger$ erscheint darin als \verb|\dagger|.

\begin{quote}\small
\begin{verbatim}
Du erstellst die Beschreibung fuer EINE Tabelle meiner Masterarbeit.

Ich gebe dir drei Dinge: (1) den vollstaendigen Fliesstext des Abschnitts, in
dem die Tabelle steht, (2) den LaTeX-Code der Tabelle selbst, (3) die
zugrundeliegende CSV aus results/tables/.

AUFGABE
Schreibe genau zwei Dinge, auf Englisch:
  A) einen \caption{}-Text, hoechstens 25 Woerter
  B) einen Notes-Block, hoechstens 90 Woerter
Bei einer Tabelle ohne Panels und ohne Marker entfaellt B ersatzlos.

WAS IN DIE NOTES GEHOERT - und sonst nichts:
  1. Definition jeder Spalte, deren Bedeutung sich nicht aus der Kopfzeile
     ergibt UND die nicht bereits in Kapitel 7 (Evaluation Framework)
     definiert ist. Ist sie dort definiert: nur \cref-Verweis, keine
     Wiederholung.
  2. Stichprobe und Einheit der Zelle (z. B. "127 monthly observations per
     country, ten countries").
  3. Verwendeter Test samt Korrektur, in einem Satz.
  4. Vorzeichenkonvention - welches Vorzeichen bedeutet welches Ergebnis.
     Das ist verpflichtend, sobald die Tabelle eine Groesse zeigt, bei der
     das nicht offensichtlich ist.
  5. Legende der Marker (z. B. was \dagger bedeutet).
  6. Provenienz, falls sie zum Nachvollziehen noetig ist (etwa welche Seeds
     welcher Zeile zugrunde liegen).

DREI HARTE VERBOTE
  * KEINE Zahl, die bereits im mitgelieferten Fliesstext steht. Pruefe jede
    Zahl gegen den Fliesstext, bevor du sie schreibst. Zahlen, die nur in der
    Tabelle stehen, sind erlaubt, wenn sie zur Definition einer Spalte noetig
    sind.
  * KEINE Interpretation, Bewertung oder Einordnung. Keine Saetze ueber das,
    was das Ergebnis bedeutet, warum es so ausfaellt oder wie stark es ist.
    Die Beschreibung sagt, wie man die Tabelle LIEST, nicht was sie ZEIGT.
  * KEINE Zahl, die du nicht in der Tabelle oder der CSV verifizieren kannst.
    Erfinde nichts, runde nichts um, rechne nichts nach.

STIL - halte dich messbar an diese Vorgaben:
  * Amerikanische Schreibung: normalized, behavior, favor, optimization.
    NICHT normalised, behaviour, favour, per cent.
  * Hoechstens ein Semikolon pro 90 Woerter.
  * Keine Geviertstriche.
  * Keine Dreierparallelen der Form "A, B, and C" mit gleichgebautem
    Nebensatz an jedem Glied.
  * Keine "not X but Y"- und "rather than"-Konstruktionen.
  * Durchschnittliche Satzlaenge um 18 Woerter, hoechstens 28.
  * Schlichte Aussagesaetze. Wo es passt, mit "We", "This" oder "The"
    beginnen.
  * Stil aus dem Fliesstext wird uebernommen.

KONVENTIONEN DER ARBEIT
  * Uebernimm die Notation des Fliesstextes unveraendert: d_norm, Winkler
    skill, admissible, package, cell, stream, gate.
  * Referenzen als \cref{}, Modellnamen in \texttt{}.
  * Die Arbeit verwendet mehrere unterschiedliche Vorzeichenkonventionen
    nebeneinander. Pruefe im Fliesstext, welche fuer DIESE Tabelle gilt, und
    uebernimm sie woertlich. Rate nicht.

AUSGABE
  1. Der fertige LaTeX-Code fuer \caption{} und, falls noetig, den
     Notes-Block.
  2. Falls dir im Fliesstext eine Definition fehlt, die du fuer die Notes
     braeuchtest: sag es, statt sie zu erfinden.
\end{verbatim}
\end{quote}

\subsection*{Betroffene Tabellen}

Caption und Notes-Block jeder Tabelle und der Abbildung wurden vollständig neu
erzeugt; die vorherigen Fassungen wurden verworfen. \emph{Bearbeitung} bezeichnet meine
Nachbearbeitung des erzeugten Textes: \emph{unverändert},
\emph{redaktionell} (nur Tippfehler und Umbrüche) oder \emph{überarbeitet}
(von mir neu formuliert).

\begin{tabular}{@{}p{6.2cm}llc@{}}
\toprule
\textbf{Tabelle / Abbildung} & \textbf{Datum} & \textbf{Bearbeitung} & \textbf{Zahlen gegen CSV geprüft} \\
\midrule
\cref{tab:features_overview}          & 21.8.2026 & unver\"andert & $\square$ \\
\cref{tab:publication-lags}           & 21.8.2026 & unver\"andert & $\square$ \\
\cref{tab:refit_blocks}               & 21.8.2026 & unver\"andert & $\square$ \\
\cref{tab:model_selection}            & 21.8.2026 & unver\"andert & $\square$ \\
\cref{tab:cp_methods}                 & 21.8.2026 & unver\"andert & $\square$ \\
\cref{tab:res_gate}                   & 21.8.2026 & unver\"andert & $\square$ \\
\cref{tab:res_point}                  & 21.8.2026 & unver\"andert & $\square$ \\
\cref{fig:skill_heatmap}            & 21.8.2026 & unver\"andert & $\square$ \\
\cref{tab:res_ranking}                & 21.8.2026 & unver\"andert & $\square$ \\
\cref{tab:res_tirex2_cp}              & 21.8.2026 & unver\"andert & $\square$ \\
\cref{tab:res_percountry}             & 21.8.2026 & unver\"andert & $\square$ \\
\cref{tab:rq3_results}                & 21.8.2026 & unver\"andert & $\square$ \\
\cref{tab:rq1_results}                & 21.8.2026 & unver\"andert & $\square$ \\
\cref{tab:res_contrasts}              & 21.8.2026 & unver\"andert & $\square$ \\
\cref{tab:rq2_results}                & 21.8.2026 & unver\"andert & $\square$ \\
\cref{tab:res_rq2_lstmG}              & 21.8.2026 & unver\"andert & $\square$ \\
\cref{tab:res_cov}                    & 21.8.2026 & unver\"andert & $\square$ \\
\cref{tab:res_xm}                     & 21.8.2026 & unver\"andert & $\square$ \\
\cref{tab:res_mcs}                    & 21.8.2026 & unver\"andert & $\square$ \\
\cref{tab:res_regime}                 & 21.8.2026 & unver\"andert & $\square$ \\
\cref{tab:res_robustness_pool}        & 21.8.2026 & unver\"andert & $\square$ \\
\cref{tab:res_robustness_ie}          & 21.8.2026 & unver\"andert & $\square$ \\
\cref{tab:res_robustness_seeds}       & 21.8.2026 & unver\"andert & $\square$ \\
\cref{tab:res_robustness_burnin}      & 21.8.2026 & unver\"andert & $\square$ \\
\cref{tab:lim_point}                  & 21.8.2026 & unver\"andert & $\square$ \\
\cref{tab:lim_dq}                     & 21.8.2026 & unver\"andert & $\square$ \\
\cref{tab:lim_amendments}             & 21.8.2026 & unver\"andert & $\square$ \\
\cref{tab:app_frozen}                 & 21.8.2026 & unver\"andert & $\square$ \\
\cref{tab:app_protocol}               & 21.8.2026 & unver\"andert & $\square$ \\
\cref{tab:app_nn_space}               & 21.8.2026 & unver\"andert & $\square$ \\
\cref{tab:app_lgbm_space}             & 21.8.2026 & unver\"andert & $\square$ \\
\cref{tab:app_frozen_nn}              & 21.8.2026 & unver\"andert & $\square$ \\
\cref{tab:app_frozen_lgbm}            & 21.8.2026 & unver\"andert & $\square$ \\
\cref{tab:app_frozen_arma}            & 21.8.2026 & unver\"andert & $\square$ \\
\cref{tab:app_country}                & 21.8.2026 & unver\"andert & $\square$ \\
\cref{tab:app_tail}                   & 21.8.2026 & unver\"andert & $\square$ \\
\cref{tab:app_pool}                   & 21.8.2026 & unver\"andert & $\square$ \\
\cref{tab:app_fragility}              & 21.8.2026 & unver\"andert & $\square$ \\
\bottomrule
\end{tabular}

\subsection*{Prompt-Fassungen}

\begin{tabular}{@{}llp{9cm}@{}}
\toprule
\textbf{Fassung} & \textbf{Datum} & \textbf{Änderung und Geltungsbereich} \\
\midrule
v1 & 21.8.2026 & Erstfassung, angewendet auf alle 37 Tabellen \\
\addlinespace
v2 & 21.8.2026 & Ergänzt die Stilregel ``Stil aus dem Fließtext wird
  übernommen''; die Ausgabe verlangt keine Kontrollliste mehr. Auf alle 37
  Tabellen erneut angewendet, davon zehn textlich geändert
  (\cref{tab:publication-lags}, \cref{tab:cp_methods}, \cref{tab:res_point},
  \cref{tab:res_ranking}, \cref{tab:res_tirex2_cp}, \cref{tab:res_rq2_lstmG},
  \cref{tab:res_cov}, \cref{tab:res_mcs}, \cref{tab:lim_point},
  \cref{tab:app_fragility}). \cref{fig:skill_heatmap} wurde direkt unter v2
  erzeugt. Der oben abgedruckte Prompt ist v2. \\
\bottomrule
\end{tabular}
